\documentclass[10pt]{article}

\usepackage[margin=1in]{geometry}
\usepackage{amsmath,amssymb,amsthm}
\usepackage{booktabs}
\usepackage{hyperref}
\usepackage{xcolor}
\usepackage{graphicx}
\usepackage{natbib}
\usepackage{microtype}

\newtheorem{definition}{Definition}
\newtheorem{theorem}{Theorem}
\newtheorem{proposition}{Proposition}
\newtheorem{conjecture}{Conjecture}

\title{Testing Compositional Faithfulness:\\Causal Interventions on Sparse Autoencoder Features\\in a Combinatorial Bijection Transformer}

\author{Robin Langer\\
\texttt{https://github.com/RaggedR/rsk-transformer}}

\date{July 2026}

\begin{document}
\maketitle

\begin{abstract}
Gauderis et al.\ recently proposed \emph{compositional interpretability},
a category-theoretic framework that formalises when a mechanistic explanation
is both faithful to a model's internal computation and compositionally
consistent. We provide the first empirical test of this framework by
performing causal feature interventions on sparse autoencoder (SAE)
decompositions of a transformer trained to invert combinatorial bijections.
We study two tasks: permutation RSK (where multiple algorithms exist) and
cylindric plane partitions (where the Burge local rule is the only known
algorithm). Zeroing individual SAE features reveals a striking contrast:
in the cylindric model, 29\% of causally important features are also
compositionally faithful (causally necessary \emph{and} correlated with
specific growth diagram faces), with the three most impactful features
all showing $>3\times$ lift on identified faces. In the permutation model,
only 2\% of causal features are compositionally faithful---the per-step
bumping path correlations reported in prior work are almost entirely
epiphenomenal. This contrast supports the Weak Principle of Explanatory
Optimism selectively: compositional faithfulness emerges when the task
constrains the model to a unique algorithm, but not when multiple
algorithmic strategies are available.
\end{abstract}

\section{Introduction}

Mechanistic interpretability aims to reverse-engineer neural networks into
human-understandable components. The standard approach---decompose, describe,
validate \citep{sharkey2025}---relies on correlational evidence: a feature
``looks like'' it encodes a concept, an attention head ``seems to'' perform
a specific computation. But correlation is not causation. A sparse autoencoder
(SAE) feature that fires on inputs matching a particular algorithmic step may
be decorative rather than functional, tracking a pattern without contributing
to the model's actual computation.

Gauderis et al.\ \citep{gauderis2026} recently proposed \emph{compositional
interpretability}, a category-theoretic framework that addresses this gap.
They define a \emph{compositional interpretation} as a pair of mappings---one
structural ($I_S$, labelling syntax), one behavioural ($I_C$, labelling
semantics)---that must \emph{commute}: $I_S = I_C \circ \llbracket \cdot
\rrbracket$. When this commutativity holds, the structural label assigned to
a component agrees with its behavioural role. They call this
\emph{compositional faithfulness}.

We provide the first empirical test of this framework. Our test bed is a
1.2M-parameter transformer trained to invert combinatorial bijections
\citep{langer2026rsk}, where the ground-truth algorithms are known exactly.
We study two tasks:

\begin{enumerate}
\item \textbf{Permutation RSK} ($n=10$, 100\% test accuracy): inverting the
Robinson--Schensted--Knuth correspondence, where two fundamentally different
algorithms exist---Schensted's sequential bumping and Fomin's parallel growth
diagrams via local rules.

\item \textbf{Cylindric plane partitions} (profile $(1,0,1,0,1,0,1,0)$,
99.98\% accuracy): inverting a cylindric growth diagram bijection where the
\emph{only} known algorithm is the Burge local rule applied face-by-face
\citep{langer2013}.
\end{enumerate}

For each task, prior work \citep{langer2026rsk} trained SAEs on the
transformer's residual stream and found features that \emph{correlate} with
specific algorithmic steps (bumping path steps for RSK, growth diagram faces
for cylindric). We extend this by asking: are these features \emph{causally
necessary}? We zero each SAE feature individually, patch the reconstructed
activation back into the model's forward pass, and measure accuracy drop.

The results reveal a sharp contrast between the two tasks, with direct
implications for the compositional interpretability framework.

\section{Background}

\subsection{Compositional Interpretability}

We briefly summarise the framework of \citet{gauderis2026}. A
\emph{compositional model} $M = (D, \mathbf{S}, \mathbf{C},
\llbracket\cdot\rrbracket)$ consists of a string diagram $D$ in a syntactic
monoidal category $\mathbf{S}$, a semantic monoidal category $\mathbf{C}$,
and a representation functor $\llbracket\cdot\rrbracket : \mathbf{S} \to
\mathbf{C}$ mapping structure to behaviour. A \emph{compositional
interpretation} is a pair of maps $I_S : \mathbf{S} \to H$ (syntactic
grounding) and $I_C : \mathbf{C} \to H$ (semantic grounding) into a space
$H$ of human-interpretable terms, satisfying the commutativity condition
$I_S = I_C \circ \llbracket\cdot\rrbracket$.

\begin{definition}[Compositional faithfulness, {\citep[Def.~3.1]{gauderis2026}}]
A compositional interpretation $I_S : \mathbf{S} \to H$ is
\emph{compositionally faithful} if there exists a reconstruction
$I_C^* : H \to \mathbf{C}$ such that for every generator $f$ of $D$,
the semantic distortion
$d(\llbracket f \rrbracket, I_C^*(I_S(f))) \leq \epsilon$.
\end{definition}

Gauderis et al.\ classify existing mechanistic interpretability methods into
a five-level \emph{refinement hierarchy} based on how the refinement functor
$R : \mathbf{S} \to \mathbf{S}'$ is constrained. SAE-based analysis operates
at Level~2 (piecewise): each component is refined in isolation, with a
per-component fidelity bound.

\begin{conjecture}[Weak Principle of Explanatory Optimism,
{\citep[Conj.~4.3]{gauderis2026}}]
Most important behaviour of an intelligence with implicit explanatory
knowledge lies within some human-understandable complexity class.
\end{conjecture}

If this conjecture holds, syntactic compression reliably surfaces
interpretable structure, because the regularities the model exploits align
with human-understandable concepts.

\subsection{The RSK Transformer}

The RSK transformer \citep{langer2026rsk} is a 1.2M-parameter encoder-only
transformer that learns to invert combinatorial bijections. Given a pair of
tableaux $(P, Q)$ encoded as structured 2D tokens (value + row + column +
tableau identity embeddings, summed into 128-d vectors), the model predicts
the original permutation $\sigma$. The same architecture handles permutation
RSK, matrix RSK, reverse plane partitions (Hillman--Grassl), and cylindric
plane partitions, achieving near-perfect accuracy on all tasks.

Prior interpretability analysis trained TopK SAEs ($d_{\text{dict}} = 1024$,
$k = 32$) on residual stream activations at each of the 6 transformer layers.
Two families of features were identified:

\begin{itemize}
\item \textbf{Permutation RSK}: \emph{insertion-order detectors} with
monotonically increasing per-step lift, and \emph{step-specific Q-entry
locators} that fire on particular $(v, r, c)$ combinations in $Q$.

\item \textbf{Cylindric}: \emph{growth diagram face detectors} that fire on
the partition triples $(i{-}1, i, i{+}1)$ involved in specific Burge local
rule applications, at $3.0$--$3.3\times$ lift.
\end{itemize}

These findings are correlational. Our experiment tests whether they are
causal.

\section{Experiment Design}

\subsection{Causal Feature Ablation}

For a trained model $M$, a trained SAE at layer $\ell$, and a feature index
$j$, we define the \emph{single-feature ablation} as follows:

\begin{enumerate}
\item Run the forward pass up to layer $\ell$.
\item Intercept the residual stream activation $h \in \mathbb{R}^{b \times t
\times d}$.
\item Encode through the SAE: normalise, apply encoder, apply TopK sparsity
to obtain sparse activations $z \in \mathbb{R}^{d_{\text{dict}}}$.
\item Set $z_j = 0$ (the intervention).
\item Decode back to $\hat{h}$ and substitute into the forward pass.
\item Measure downstream accuracy.
\end{enumerate}

The \emph{causal importance} of feature $j$ is $\Delta_j = a_{\text{SAE}} -
a_j$, where $a_{\text{SAE}}$ is the accuracy with full SAE reconstruction
(no features zeroed) and $a_j$ is the accuracy with feature $j$ zeroed.
Using the SAE baseline rather than the original model isolates the effect of
the single feature from the SAE reconstruction error.

\subsection{Cross-Referencing Causal and Correlational Evidence}

For each feature, we have two independent measurements:

\begin{itemize}
\item \textbf{Correlational}: the \emph{per-step lift} (permutation) or
\emph{per-face lift} (cylindric)---how much more likely the feature is to
fire on tokens involved in a specific algorithmic step compared to its base
rate.

\item \textbf{Causal}: $\Delta_j$, the accuracy drop from zeroing the feature.
\end{itemize}

A feature is \emph{compositionally faithful} (in the sense of
\citet{gauderis2026}) if it scores high on \emph{both} measures: the
structural label assigned by the correlational analysis ($I_S$) agrees with
the feature's functional role as revealed by the causal intervention ($I_C$).
We operationalise this as $\Delta_j > 0.1\%$ (causal) and lift $> 1.5\times$
(correlated).

\subsection{Setup}

\begin{itemize}
\item \textbf{Permutation RSK}: model \texttt{encoder\_n10} (100\% greedy
exact match on 50K test samples), SAEs from
\texttt{sae\_encoder\_n10}, layer~2 (the critical computation layer; ablating
it drops accuracy to 62\%).

\item \textbf{Cylindric}: model \texttt{encoder\_cyl\_10101010\_m3} (99.98\%
exact match), SAEs from \texttt{sae\_cyl\_10101010}, layer~2.

\item \textbf{Evaluation}: 10,000 test samples (permutation) and 1,000 test
samples (cylindric), seeded for reproducibility.

\item \textbf{Sanity checks}: zeroing all features drops accuracy to 0\%
(confirming features carry real information); zeroing dead features produces
$\Delta \approx 0$.
\end{itemize}

\section{Results}

\subsection{Permutation RSK}

Table~\ref{tab:perm} shows the top 15 features by causal importance at
layer~2.

\begin{table}[h]
\centering
\caption{Permutation RSK (layer~2, $n{=}10$, 10K test samples). SAE
baseline: 100.00\% greedy exact match. No feature in the top 30 by causal
importance has lift $> 1.5$.}
\label{tab:perm}
\small
\begin{tabular}{rccccl}
\toprule
Feature & $\Delta$ Exact & $\Delta$ Per-pos & Lift & Best Step & Assessment \\
\midrule
664 & 15.34\% & 3.07\% & 1.11 & $k{=}1$ & causal only \\
895 & 12.58\% & 3.03\% & 1.15 & $k{=}9$ & causal only \\
361 & 10.67\% & 2.87\% & 1.39 & $k{=}1$ & causal only \\
387 & 10.49\% & 2.91\% & 1.49 & $k{=}0$ & causal only \\
678 & 10.25\% & 2.21\% & 1.15 & $k{=}7$ & causal only \\
657 &  8.16\% & 1.67\% & 1.08 & $k{=}7$ & causal only \\
230 &  7.32\% & 1.94\% & 1.08 & $k{=}4$ & causal only \\
406 &  6.47\% & 1.61\% & 1.03 & $k{=}2$ & causal only \\
409 &  5.37\% & 1.46\% & 1.03 & $k{=}3$ & causal only \\
 23 &  4.81\% & 1.21\% & 1.11 & $k{=}7$ & causal only \\
564 &  3.94\% & 1.34\% & 1.10 & $k{=}8$ & causal only \\
129 &  3.54\% & 0.88\% & 1.02 & $k{=}3$ & causal only \\
611 &  2.68\% & 0.71\% & 1.13 & $k{=}8$ & causal only \\
753 &  2.43\% & 0.49\% & 1.26 & $k{=}8$ & causal only \\
715 &  1.80\% & 0.40\% & 1.19 & $k{=}1$ & causal only \\
\bottomrule
\end{tabular}
\end{table}

Of 405 alive features:
\begin{itemize}
\item 43 are causally important ($\Delta > 0.1\%$)
\item 226 are correlationally strong (lift $> 1.5\times$)
\item \textbf{1} is both (2\%)
\item Top-20 overlap between causal and correlational rankings: \textbf{0}
features
\end{itemize}

The per-step bumping path correlations are almost entirely epiphenomenal. The
most causally important feature (\#664, $\Delta = 15.3\%$) has lift 1.11 on
step $k{=}1$---it fires on step~1's bumping path barely more than on any
other P-token. Meanwhile, the 226 features with lift $> 1.5\times$ are
individually dispensable: the model routes around any single one.

\subsection{Cylindric Plane Partitions}

Table~\ref{tab:cyl} shows the top 15 features by causal importance at
layer~2.

\begin{table}[h]
\centering
\caption{Cylindric plane partitions (layer~2, profile $(1,0,1,0,1,0,1,0)$,
1K test samples). SAE baseline: 99.90\% exact match. Three of the top four
causal features have lift $> 3\times$.}
\label{tab:cyl}
\small
\begin{tabular}{rccccl}
\toprule
Feature & $\Delta$ Exact & $\Delta$ Per-pos & Lift & Best Face & Assessment \\
\midrule
861 & 38.10\% & 3.81\% & 3.16 & $f{=}0$ & \textbf{faithful} $\star$ \\
1006 & 31.90\% & 3.21\% & 3.28 & $f{=}6$ & \textbf{faithful} $\star$ \\
469 & 23.50\% & 4.56\% & 1.11 & $f{=}3$ & causal only \\
177 & 22.90\% & 4.14\% & 1.06 & $f{=}6$ & causal only \\
 38 & 21.90\% & 4.23\% & 1.04 & $f{=}4$ & causal only \\
279 & 19.80\% & 3.94\% & 1.23 & $f{=}2$ & causal only \\
185 & 18.00\% & 3.43\% & 1.16 & $f{=}6$ & causal only \\
294 & 15.40\% & 1.54\% & 3.20 & $f{=}6$ & \textbf{faithful} $\star$ \\
129 & 14.70\% & 1.91\% & 1.17 & $f{=}1$ & causal only \\
746 & 13.90\% & 2.23\% & 1.08 & $f{=}0$ & causal only \\
 91 & 13.30\% & 1.36\% & 1.91 & $f{=}3$ & \textbf{faithful} $\star$ \\
865 &  9.30\% & 1.99\% & 1.11 & $f{=}4$ & causal only \\
230 &  9.00\% & 1.03\% & 1.24 & $f{=}6$ & causal only \\
516 &  8.40\% & 0.89\% & 1.28 & $f{=}2$ & causal only \\
551 &  5.20\% & 0.54\% & 1.31 & $f{=}1$ & causal only \\
\bottomrule
\end{tabular}
\end{table}

Of 356 alive features:
\begin{itemize}
\item 49 are causally important ($\Delta > 0.1\%$)
\item 267 are correlationally strong (lift $> 1.5\times$)
\item \textbf{14} are both (29\%)
\item The three most causally impactful features (\#861, \#1006, \#294) all
have lift $> 3\times$ on specific growth diagram faces
\end{itemize}

\subsection{Comparison}

\begin{table}[h]
\centering
\caption{Summary comparison: compositional faithfulness by task.}
\label{tab:comparison}
\begin{tabular}{lcc}
\toprule
& Permutation RSK & Cylindric \\
\midrule
Causal features ($\Delta > 0.1\%$) & 43 & 49 \\
Correlated features (lift $> 1.5$) & 226 & 267 \\
Compositionally faithful & 1 (2\%) & 14 (29\%) \\
Top-20 overlap & 0 & 1 \\
Top causal feature: $\Delta$ & 15.3\% & 38.1\% \\
Top causal feature: lift & 1.11 & 3.16 \\
\bottomrule
\end{tabular}
\end{table}

\section{Discussion}

\subsection{Task Constraints Determine Compositional Faithfulness}

The central finding is that compositional faithfulness---the alignment
between correlational structure and causal importance---depends on the
\emph{task}, not the model architecture.

The permutation RSK task admits multiple algorithms: Schensted's sequential
bumping and Fomin's parallel growth diagrams. The transformer is free to
learn any strategy that inverts the bijection, and the SAE decomposition
reveals that it has learned a \emph{distributed} representation where no
single feature ``owns'' a specific algorithmic step. The per-step bumping
path correlations, while statistically real (many features have lift
$> 1.5\times$), are epiphenomenal---the model can compute the correct
output without any individual correlated feature.

The cylindric task admits only one algorithm: the Burge local rule applied
face-by-face through a growth diagram \citep{langer2013}. The computation
\emph{must} be localised to faces, so SAE features that correlate with
specific faces are also causally necessary. The three most impactful features
all have lift $> 3\times$ on identified faces, and zeroing any one of them
drops accuracy by 15--38\%.

In the language of \citet{gauderis2026}, the cylindric task constrains the
\emph{refinement space}: there are fewer valid decompositions of the
computation, so the model's actual syntax is closer to the unique
interpretable decomposition. The permutation task leaves the refinement
space unconstrained, and the model's chosen syntax does not align with
either of the two known human-interpretable algorithms at the
single-feature level.

\subsection{Implications for Explanatory Optimism}

The Weak Principle of Explanatory Optimism (Conjecture~4.3 of
\citealt{gauderis2026}) holds that models with implicit explanatory knowledge
will tend to represent it in human-understandable form. Our results suggest a
refinement: explanatory optimism is \emph{task-conditional}.

When the task structure forces a unique algorithmic solution, the model's
internal representation aligns with human-understandable concepts (growth
diagram faces), and SAE features are compositionally faithful. When multiple
algorithmic strategies are available, the model finds a distributed
representation that is computationally effective but not aligned with any
single human-understandable algorithm at the piecewise (L2) decomposition
level.

This does not refute explanatory optimism outright. The permutation model
may still admit a compositionally faithful interpretation at a different
refinement level---for instance, Level~3 (global, cross-layer circuits) might
reveal that \emph{groups} of features jointly implement bumping steps, even
though no single feature does so alone. Our experiment tests L2 faithfulness
only.

\subsection{Methodological Implications}

Our results carry a warning for the mechanistic interpretability community:
\textbf{correlational SAE feature analysis can be systematically misleading}.
In the permutation model, 226 features pass a reasonable correlational
threshold (lift $> 1.5\times$), yet essentially none of them are causally
important. A research programme that reports only correlational evidence
would conclude that the model has learned Schensted-like bumping steps; the
causal evidence shows this conclusion is unfounded.

The fix is straightforward: single-feature ablation is computationally cheap
(one forward pass per feature) and should be standard practice alongside
correlational analysis. The two measurements are complementary---correlation
identifies \emph{what} a feature tracks, causation identifies \emph{whether}
it matters---and only their intersection constitutes compositional
faithfulness.

\subsection{Limitations}

\begin{itemize}
\item Our analysis is restricted to Level~2 (piecewise, per-layer)
refinement. Cross-layer circuits, which may carry the permutation model's
true algorithmic structure, require Level~3 analysis.

\item The cylindric experiment uses 1,000 test samples (vs 10,000 for
permutation), introducing more noise in the $\Delta$ estimates.

\item We test only layer~2. Other layers may show different patterns of
faithfulness.

\item The lift threshold of $1.5\times$ is somewhat arbitrary. However, the
qualitative result is robust: even at a threshold of $1.2\times$, only 5 of
43 causal features qualify for the permutation model.

\item We ablate single features only. Multi-feature ablation (testing whether
groups of correlated features are jointly causal) is a natural next step.
\end{itemize}

\section{Related Work}

\textbf{Compositional interpretability.}
\citet{gauderis2026} introduce the framework we test.
\citet{tull2024} develop the categorical foundations for compositional XAI.
\citet{barbiero2025} propose defining interpretability through symmetry
invariance.

\textbf{Causal methods in interpretability.}
\citet{geiger2025} frame mechanistic interpretability as causal abstraction,
using top-down hypothesis testing to find alignments between models and
high-level explanations. Our approach is bottom-up: we start from SAE
features and test their causal necessity.
\citet{joshi2026} argue that causal evaluation is essential for
interpretability claims to generalise.

\textbf{SAE interpretability.}
\citet{bricken2023} introduce dictionary learning for transformer
decomposition.
\citet{templeton2024} scale monosemantic feature extraction to Claude~3
Sonnet. \citet{braun2025} connect SAE explanations to minimum description
length.

\textbf{ML for combinatorics.}
\citet{kvinge2025} establish the ML4AlgComb benchmark for machine learning
on algebraic combinatorics. \citet{langer2026rsk} train the RSK transformer
we study here.

\section{Conclusion}

We provide the first empirical test of the compositional interpretability
framework of \citet{gauderis2026}, using causal feature ablation on SAE
decompositions of a combinatorial bijection transformer. The central finding
is a task-dependent contrast:

\begin{itemize}
\item When the task admits a unique algorithm (cylindric plane partitions),
SAE features that correlate with algorithmic steps are also causally
necessary. Compositional faithfulness holds: 29\% of causal features are
also the strongest face-correlated features, and the top three causal
features all show $>3\times$ lift.

\item When multiple algorithms exist (permutation RSK), correlational
evidence is systematically misleading. The 226 features with strong per-step
bumping path correlations are individually dispensable; the 43 causally
necessary features show no step-specific structure.
\end{itemize}

This suggests that explanatory optimism---the conjecture that syntactic
compression surfaces human-aligned explanations---is task-conditional. It
holds when the task constrains the refinement space to a single
decomposition, but fails when the model is free to distribute computation
across features in ways that don't map onto known algorithms.

The practical implication is clear: correlational SAE analysis should always
be accompanied by causal intervention. Only the intersection of the two
constitutes a compositionally faithful interpretation.

\paragraph{Code and models.}
Source code: \url{https://github.com/RaggedR/rsk-transformer}.
Trained models: \url{https://huggingface.co/RobBobin/rsk-transformer}.

\bibliographystyle{plainnat}

\begin{thebibliography}{99}

\bibitem[Barbiero et al.(2025)]{barbiero2025}
P.~Barbiero, M.~E.~Zarlenga, F.~Termine, A.~Jamnik, M.~Marra.
\newblock Foundations of Interpretable Models.
\newblock August 2025.

\bibitem[Braun et al.(2025)]{braun2025}
D.~Braun, L.~Bushnaq, L.~Heimersheim, S.~Mendel, J.~Sharkey, L.
\newblock Interpretability in parameter space: Minimizing mechanistic
description length with attribution-based parameter decomposition.
\newblock \emph{arXiv:2501.14926}, 2025.

\bibitem[Bricken et al.(2023)]{bricken2023}
T.~Bricken, A.~Templeton, J.~Batson, et al.
\newblock Towards monosemanticity: Decomposing language models with dictionary
learning.
\newblock \emph{Transformer Circuits Thread}, 2023.

\bibitem[Gauderis et al.(2026)]{gauderis2026}
W.~Gauderis, T.~Dooms, S.~T.~Homer, K.~Ayonrinde, G.~A.~Wiggins.
\newblock From mechanistic to compositional interpretability.
\newblock In \emph{2nd Workshop on Compositional Learning, ICML 2026}, 2026.
\newblock \emph{arXiv:2605.08934}.

\bibitem[Geiger et al.(2025)]{geiger2025}
A.~Geiger, D.~Ibeling, D.~Zur, A.~Chaudhary, et al.
\newblock Causal abstraction: A theoretical foundation for mechanistic
interpretability.
\newblock May 2025.

\bibitem[Joshi et al.(2026)]{joshi2026}
S.~Joshi, S.~Mueller, A.~Klindt, D.~Brendel, W.~Reizinger, P.~Sridhar, D.
\newblock Causality is key for interpretability claims to generalise.
\newblock \emph{arXiv:2602.16698}, February 2026.

\bibitem[Kvinge et al.(2025)]{kvinge2025}
H.~Kvinge, D.~Brown, H.~Jenne, R.~Jasper, S.~Scullen.
\newblock ML4AlgComb: Machine learning for algebraic combinatorics.
\newblock \emph{arXiv:2503.06366}, 2025.

\bibitem[Langer(2013)]{langer2013}
R.~Langer.
\newblock \emph{Cylindric plane partitions, Lambda determinants, Commutants in
semicircular systems}.
\newblock PhD thesis, Universit\'{e} Paris-Est, 2013.
\newblock \emph{arXiv:2110.12629}.

\bibitem[Langer(2026)]{langer2026rsk}
R.~Langer.
\newblock Learning combinatorial bijections with transformers: Inverse RSK,
growth diagrams, and interpretable features.
\newblock April 2026.

\bibitem[Sharkey et al.(2025)]{sharkey2025}
L.~Sharkey, B.~Chughtai, J.~Batson, et al.
\newblock Open problems in mechanistic interpretability.
\newblock January 2025.

\bibitem[Templeton et al.(2024)]{templeton2024}
A.~Templeton, T.~Conerly, J.~Marcus, et al.
\newblock Scaling monosemanticity: Extracting interpretable features from
Claude~3 Sonnet.
\newblock \emph{Transformer Circuits Thread}, 2024.

\bibitem[Tull et al.(2024)]{tull2024}
S.~Tull, R.~Lorenz, R.~Clark, S.~Khan, I.~Coecke, B.
\newblock Towards compositional interpretability for XAI.
\newblock June 2024.

\end{thebibliography}

\end{document}
